\chapter{Tardive dyskinesia}
\index{Tardive dyskinesia}

\medskip
\noindent\textit{Educational reference; always seek professional advice for individual care.}

\section*{Definition and Overview}
Tardive dyskinesia is a movement disorder that develops after long-term treatment with medicines that block dopamine receptors in the brain. The word "tardive" means late, and the condition appears only after months or years of exposure, in contrast to the early movement problems some people experience soon after starting these medicines. It is characterised by involuntary, repetitive, often twisting movements, most commonly affecting the face, mouth, and tongue, but also involving the arms, legs, and trunk. Tardive dyskinesia is caused mainly by antipsychotic medicines, used to treat schizophrenia, bipolar disorder, and other conditions, but it can also be caused by certain medicines used for nausea and vertigo, such as metoclopramide and prochlorperazine. It is important because it can be distressing and stigmatising, can interfere with eating, speaking, and daily life, and may be permanent in some people even after the offending medicine is stopped. For people who need antipsychotics to control serious mental illness, the aim is not simply to stop treatment, but to find the lowest effective dose or an alternative medicine that is less likely to cause the problem, and to recognise the earliest signs. This chapter explains how tardive dyskinesia develops, who is at risk, how it is recognised and distinguished from other movement disorders, and the modern treatments available.

\section*{Epidemiology}
Tardive dyskinesia is one of the most important long-term side effects of antipsychotic treatment. In people taking older ("typical") antipsychotics such as haloperidol, it is estimated that roughly five percent develop tardive dyskinesia each year of treatment, so that after several years a substantial minority are affected. The newer "atypical" antipsychotics such as risperidone, olanzapine, and quetiapine carry a lower but still real risk, particularly at higher doses. Older age is the strongest risk factor, and women are affected somewhat more often than men. The condition is far less common in younger people, but it does occur, and it can affect children treated with anti-nausea medicines. Because antipsychotics are used widely and for long periods, tardive dyskinesia affects a meaningful number of people in many countries, and it is also seen in people who have taken metoclopramide or prochlorperazine for months for nausea or migraine. Rates have fallen somewhat since the shift to atypical antipsychotics, but the condition has not disappeared, and careful monitoring remains an essential part of prescribing these medicines.

\section*{Aetiology and Pathophysiology}
Tardive dyskinesia is caused by medicines that block dopamine receptors, but the exact mechanism by which long-term blockade produces abnormal movements is still not fully understood. The leading explanation involves changes in the sensitivity of dopamine pathways in a region of the brain called the basal ganglia, which controls fine movement.

\begin{itemize}
    \item \textbf{Antipsychotics:} the main cause; typical antipsychotics such as haloperidol, fluphenazine, and trifluoperazine carry the highest risk, while atypical antipsychotics such as risperidone, olanzapine, quetiapine, aripiprazole, and clozapine carry lower risk at lower doses
    \item \textbf{Anti-nausea and vertigo medicines:} metoclopramide and prochlorperazine block dopamine and cause the condition when used for months or years; domperidone is a similar medicine with lower brain penetration
    \item \textbf{Dopamine receptor supersensitivity:} long-term blockade is thought to cause the receptors to become abnormally sensitive, so that normal dopamine signalling produces exaggerated, uncontrolled movements
    \item \textbf{Oxidative stress:} dopamine metabolism generates damaging free radicals, and cumulative oxidative injury to the basal ganglia may contribute, especially in older people
    \item \textbf{Genetic susceptibility:} not everyone exposed develops the condition, and variation in genes handling dopamine metabolism or oxidative protection is thought to explain some of this difference
    \item \textbf{Other medicines:} rarely, some antidepressants, lithium, or anticonvulsants have been associated with dyskinetic movements, though the link is far weaker than for dopamine-blocking drugs
\end{itemize}

The practical implication is that risk is cumulative: the higher the dose and the longer the treatment, the greater the chance of tardive dyskinesia. This is why doctors review dose and need regularly and look for the earliest signs before movements become fixed.

\section*{Clinical Features}
Tardive dyskinesia is diagnosed on the movements themselves, which are involuntary, repetitive, and often stereotyped. The features vary from person to person, but commonly include:

\begin{itemize}
    \item \textbf{Lip smacking and puckering:} repetitive, involuntary movements of the lips that can resemble chewing or kissing
    \item \textbf{Tongue movements:} the tongue may dart in and out, twist, or push against the inside of the cheek, and this is often one of the earliest signs
    \item \textbf{Grimacing:} uncontrolled pulling of the facial muscles into expressions of discomfort
    \item \textbf{Rapid blinking or eye rolling:} the eyelids and eyes may move involuntarily
    \item \textbf{Jaw movements:} chewing, clamping, or sideways grinding motions of the jaw, which can damage teeth
    \item \textbf{Choreiform limb movements:} quick, jerky, dance-like movements of the fingers, hands, or feet, sometimes looking fidgety
    \item \textbf{Athetoid movements:} slower, writhing, twisting movements of the arms, legs, or fingers
    \item \textbf{Truncal movements:} rocking, swaying, or arching of the hips, pelvis, or shoulders when sitting or standing
    \item \textbf{Respiratory involvement:} irregular grunting, sighing, or breathing in some people with more severe disease
    \item \textbf{Suppression and worsening:} the movements reduce when the person relaxes and concentrates and increase with stress or anxiety; they stop in sleep
\end{itemize}

The movements are usually mild at first and may be barely noticeable to others, which is why regular examination matters: the earliest signs, such as subtle tongue movements, are the easiest to reverse.

\section*{Differential Diagnosis}
Not every involuntary movement in someone taking an antipsychotic is tardive dyskinesia, and the distinction matters because treatment differs.

\begin{itemize}
    \item \textbf{Drug-induced parkinsonism:} rigidity, shuffling, and tremor appear early in treatment and usually improve when the dose is reduced, whereas tardive dyskinesia appears late and involves rapid, writhing movements
    \item \textbf{Akathisia:} a feeling of inner restlessness with an urge to pace, which appears early and is not the same as the repetitive involuntary movements of tardive dyskinesia
    \item \textbf{Acute dystonia:} sudden muscle spasms, often of the neck or eyes, occurring within days of starting a dopamine-blocking medicine, and treated urgently
    \item \textbf{Huntington's disease:} an inherited condition with chorea and cognitive change, distinguished by family history, age of onset, and genetic testing
    \item \textbf{Tics and Tourette syndrome:} tics are semi-voluntary and preceded by a premonitory urge, whereas tardive dyskinesia usually starts later in life
    \item \textbf{Stereotypies:} repetitive, apparently purposeful movements seen in autism and other neurodevelopmental conditions, not related to medicines
    \item \textbf{Spontaneous orofacial dyskinesia:} in older people, especially those with missing teeth or poor dentures, similar movements can occur without medicine exposure
    \item \textbf{Tremor:} parkinsonian and essential tremors are rhythmic, while tardive dyskinesia movements are irregular and non-rhythmic
\end{itemize}

\section*{When to Seek Medical Care}
Anyone taking a dopamine-blocking medicine who notices new, involuntary movements of the face, mouth, tongue, hands, or trunk should see their doctor promptly. Do not stop the medicine yourself, because sudden withdrawal can itself cause severe movement problems and can destabilise the mental illness being treated. The doctor will review the dose and the need for the medicine and discuss options such as switching to a lower-risk medicine or adding treatment. Family members and carers should also speak up, because the person affected is often unaware of their own movements.

\textbf{Contact your local emergency services for:}
\begin{itemize}
    \item Difficulty breathing, or noisy, irregular breathing that is new
    \item Choking or coughing when eating or drinking, or inability to swallow saliva
    \item A sudden crisis with stiffening or sustained muscle spasms (dystonia), fever, and confusion
\end{itemize}

\section*{Diagnosis and Investigations}
Tardive dyskinesia is a clinical diagnosis: there is no blood test that confirms it, and the diagnosis rests on the history of medicine exposure and the pattern of movements.

\begin{itemize}
    \item \textbf{Drug history:} the most important step is documenting which dopamine-blocking medicines were taken, at what doses, and for how long, including anti-nausea medicines easily forgotten
    \item \textbf{Clinical examination:} the doctor observes the movements at rest, while the person is distracted, and while performing other tasks, looking for the typical orofacial and limb patterns
    \item \textbf{Standardised rating scales:} tools such as the Abnormal Involuntary Movement Scale (AIMS) score the severity in different body areas and track change over time
    \item \textbf{Exclusion of other causes:} a neurological review may be arranged to rule out Parkinson's disease, Huntington's disease, or thyroid problems where the picture is unclear
    \item \textbf{Blood tests and imaging:} thyroid function, copper studies, or a brain scan are ordered only if a systemic or structural cause is suspected
\end{itemize}

\section*{Management}
The management of tardive dyskinesia has improved greatly in recent years, and there is much that can be done even when the movements have been present for some time.

\begin{itemize}
    \item \textbf{Early detection and review:} recognising the first signs and reviewing the medicine is the most important step, because early movements are more likely to reverse
    \item \textbf{Dose reduction or withdrawal:} where the underlying condition allows, the dose is reduced or stopped gradually, under medical supervision
    \item \textbf{Switching medicines:} when antipsychotic treatment must continue, a switch to a lower-risk atypical agent such as quetiapine or clozapine is arranged with the treating psychiatrist
    \item \textbf{VMAT2 inhibitors:} medicines such as deutetrabenazine and valbenazine reduce the release of dopamine and are now approved treatments that reduce the severity of movements in many people; some are covered in some public systems for eligible patients
    \item \textbf{Botulinum toxin injections:} injections into the affected muscles can help people with bothersome focal movements, such as eyelid or neck involvement
    \item \textbf{Other medicines:} clonazepam, amantadine, or certain other agents have been used with modest benefit in some people
    \item \textbf{Medication review for antiemetics:} people taking metoclopramide or prochlorperazine long term should discuss stopping or switching, as these are avoidable causes
    \item \textbf{Multidisciplinary care:} physiotherapy, dental review for tooth damage, and swallowing assessment if eating is affected all form part of comprehensive care
\end{itemize}

\section*{Complications}
Tardive dyskinesia is more than a cosmetic problem, and its complications can be serious.

\begin{itemize}
    \item \textbf{Social and psychological distress:} visible movements can cause embarrassment, stigmatisation, social withdrawal, and worsening of mood
    \item \textbf{Dental damage:} jaw and tongue movements can grind down teeth and irritate the gums
    \item \textbf{Difficulty eating and weight loss:} uncontrolled mouth and jaw movements can make chewing and swallowing hard, leading to poor nutrition
    \item \textbf{Speech problems:} tongue and respiratory movements can make speech unclear or interrupted
    \item \textbf{Respiratory complications:} in severe cases, irregular breathing or aspiration of food into the lungs can occur
    \item \textbf{Persistence:} a significant minority have permanent movements even after the offending medicine is stopped
\end{itemize}

\section*{Prognosis and Outcomes}
The outlook for tardive dyskinesia depends mainly on how quickly it is recognised and how flexible the treatment can be. When movements are detected early and the medicine can be reduced or changed, many people improve substantially, and some recover completely within weeks to months. In people who must continue high-dose antipsychotics, the movements may persist or worsen, and a smaller group are left with permanent dyskinesia after the medicine is stopped. The modern VMAT2 inhibitor treatments offer real benefit: clinical trials show they reduce the severity of movements by a meaningful margin in a large proportion of people, giving hope where there was little before. Even in people with persistent movements, supportive measures such as dental care, swallowing advice, and psychological support make the condition manageable, and the underlying psychiatric illness can continue to be treated effectively. The message for patients and families is to stay alert to early signs, keep every medication review, and ask specifically about movement side effects, because the earlier the problem is addressed, the better the outcome.

\section*{Prevention and Self-Care}
\begin{itemize}
    \item Use the lowest effective dose of any antipsychotic or anti-nausea medicine, and ask your doctor to review the need regularly
    \item Ask about switching to a lower-risk atypical antipsychotic if you take a typical antipsychotic long term
    \item Do not take metoclopramide or prochlorperazine for longer than a few days without discussing it with your doctor
    \item Attend every scheduled review, and ask a family member or carer to tell you if they notice any new movements
    \item Keep a current medication list and show it to every doctor, dentist, and pharmacist
    \item If you notice any new movements, report them early rather than waiting for the next appointment
    \item Never stop or reduce antipsychotics without medical advice, as sudden withdrawal can cause severe problems
    \item Maintain good dental care and report any tooth grinding or soreness promptly
\end{itemize}

\section*{Sources and Further Reading}
The following organisations provide reliable information about tardive dyskinesia and the safe use of antipsychotic medicines.

\begin{itemize}
    \item healthdirect Australia. \textit{Tardive dyskinesia: causes, symptoms, and treatment.} https://www.healthdirect.gov.au/
    \item Mayo Clinic. \textit{Tardive dyskinesia: symptoms and treatment options.} https://www.mayoclinic.org/
    \item NHS. \textit{Antipsychotic medicines: side effects and monitoring.} https://www.nhs.uk/
    \item NPS MedicineWise. \textit{Antipsychotics and movement side effects.} https://www.nps.org.au/
    \item Therapeutic Goods Administration. \textit{Medicine information for VMAT2 inhibitors.} https://www.tga.gov.au/
    \item Australian Department of Health and Aged Care. \textit{Mental health treatment and medicines resources.} https://www.health.gov.au/
\end{itemize}

\clearpage
